\Subhead{3.1.1}{Случайные величины и процессы, способы их описания. Стационарный случайный процесс}

Случайная величина $X$ задаётся функцией распределения $F_X(x)=\Pr(X\leq x)$; при наличии плотности $p_X=F'_X$ среднее $\langle g(X)\rangle=\int g(x)p_X(x)\dd x$. Случайный процесс $X(t)$ — семейство случайных величин: один опыт даёт реализацию, а повторение опыта — ансамбль. Полное описание задаёт согласованная система многомерных распределений $p(x_1,t_1;\ldots;x_n,t_n)$.

Строгая стационарность означает инвариантность всех распределений относительно общего сдвига времени. Стационарность в широком смысле требует конечного постоянного среднего $m$ и корреляции, зависящей только от разности времён:
\begin{equation}
 C_X(\tau)=\langle[X(t+\tau)-m][X(t)-m]\rangle.
\end{equation}
Она достаточна для линейной спектральной теории второго порядка, но не определяет высшие моменты. У гауссовского процесса среднее и ковариация определяют все распределения, поэтому широкая стационарность влечёт строгую. Момент второго порядка $R_X(\tau)=C_X(\tau)+m^2$ нужно отличать от ковариации: постоянное среднее создаёт спектральную дельта-линию при нулевой частоте.
